% ======================================================================
% Paper 3 Technical Supplement v1.0
% Companion to Paper 3 --- Ledgers: Admissible Dynamics, Entropy, and the
% Arrow of Time.
%
% v1.0 (2026-04-18): Initial Technical Supplement, built in response to
%   the Paper 3 lit-review audit (2026-04-18).  Mirrors the structure of
%   the Paper 1 / Paper 2 / Paper 7 Technical Supplements.  Six formal
%   sections cover the main theorems at proof-dense scope; a red-team
%   section flags non-A1/non-PLEC hypotheses at point of use; theorem
%   index and dependency appendix close the document.
% ======================================================================

\documentclass[11pt]{article}

\usepackage[T1]{fontenc}
\usepackage[utf8]{inputenc}
\usepackage[a4paper,margin=1in]{geometry}
\usepackage{setspace}
\usepackage{parskip}
\usepackage{enumitem}
\usepackage{booktabs}
\usepackage{array}
\usepackage{caption}
\usepackage{amsmath,amssymb,amsfonts,amsthm}
\usepackage{mathtools}
\usepackage[most]{tcolorbox}
\usepackage{titlesec}
\usepackage{fancyhdr}
\usepackage[colorlinks=true,linkcolor=black,citecolor=black,urlcolor=black]{hyperref}

\newtheorem{theorem}{Theorem}[section]
\newtheorem{lemma}[theorem]{Lemma}
\newtheorem{proposition}[theorem]{Proposition}
\newtheorem{corollary}[theorem]{Corollary}
\newtheorem{definition}[theorem]{Definition}
\theoremstyle{remark}
\newtheorem{remark}[theorem]{Remark}

\titleformat{\section}{\Large\bfseries}{\thesection.}{0.5em}{}
\titleformat{\subsection}{\large\bfseries}{\thesubsection}{0.5em}{}

\pagestyle{fancy}
\fancyhf{}
\fancyhead[L]{\small\itshape Paper 3 Technical Supplement}
\fancyhead[R]{\small\itshape E.\,Brooke}
\fancyfoot[C]{\thepage}
\renewcommand{\headrulewidth}{0.4pt}

\tcbuselibrary{skins,breakable}

\tcbset{
    apfbox/.style={
        enhanced,
        colback=black!3,
        frame hidden,
        borderline={0.5pt}{0pt}{black},
        arc=0pt,
        left=8pt, right=8pt, top=8pt, bottom=8pt,
        fonttitle=\bfseries\color{black},
        coltitle=black,
        detach title,
        before upper={\tcbtitle\par\smallskip},
        before skip=14pt,
        after skip=14pt,
        sharp corners,
        breakable,
    }
}

\tcbset{
    redteambox/.style={
        enhanced,
        colback=black!5,
        frame hidden,
        borderline={0.5pt}{0pt}{black!70},
        arc=0pt,
        left=8pt, right=8pt, top=8pt, bottom=8pt,
        fonttitle=\bfseries\color{black!80},
        coltitle=black!80,
        detach title,
        before upper={\tcbtitle\par\smallskip},
        before skip=14pt,
        after skip=14pt,
        sharp corners,
        breakable,
    }
}

\newcommand{\coderef}[2]{\smallskip\noindent\textit{Code reference:} \texttt{#1} in \texttt{#2}.}
\newcommand{\epitag}[1]{\textsf{\footnotesize[#1]}}

\setstretch{1.05}
\setlist[itemize]{itemsep=3pt,topsep=4pt}
\setlist[enumerate]{itemsep=3pt,topsep=4pt}

\title{\textbf{Paper 3 Technical Supplement}\\[0.4em]
\large Admissible Dynamics, Entropy, and the Arrow of Time: Formal Proofs\\[0.2em]
\normalsize Companion to Paper~3 of the Admissibility Physics Framework}
\author{E.\ S.\ Brooke}
\date{Version 1.0 --- April 2026}

\begin{document}

\maketitle

\noindent\textit{Independent Researcher, San Anselmo, California, USA}\\
\textit{Corresponding author: brooke.ethan@gmail.com}\\
\textit{ORCID: \href{https://orcid.org/0009-0001-2261-4682}{0009-0001-2261-4682}}

\bigskip

\begin{abstract}
\noindent This Technical Supplement provides proof-dense formal treatment of
the main theorems of Paper~3 --- \emph{Ledgers: Admissible Dynamics,
Entropy, and the Arrow of Time}.  The main paper states the theorems with
proof sketches; the Supplement carries the full formal burden.  The
organization mirrors the Paper~1, Paper~2, and Paper~7 Technical Supplements.
Six formal sections cover, in order: complete positivity as a theorem
($T_{\mathrm{CPTP}}$) in explicit comparison to the Stinespring--Choi--GKSL
infrastructure~\cite{Stinespring1955,Choi1975,Lindblad1976,GKS1976,Kraus1983};
the directed-enforcement multiplier with $\kappa=2$ derived from
irreversibility and locality ($T_\kappa$); von Neumann entropy as committed
capacity ($T_{\mathrm{entropy}}$) anchored to the strong-subadditivity
backbone~\cite{LiebRuskai1973,Wehrl1978}; the second law and the arrow of
time as structural consequences of $L_{\mathrm{irr}}$; admissibility carrying
capacity (ACC) and horizon entropy $S=A/4$ with proof of the $1/2\times 1/2$
bookkeeping; and the Fisher-information metric on admissible states with RG
flow as Fisher gradient descent~\cite{BraunsteinCaves1994,Petz1996,AmariNagaoka2000}.
A red-team section flags three additional hypotheses at point of use
(H1 Markovian-composition closure; H2 Hamiltonian-constraint redundancy;
H3 constraint-functional monotonicity).  Theorem index and dependency-diagram
appendices close the document.
\end{abstract}

\bigskip
\tableofcontents
\bigskip


% ======================================================================
\section{Front matter: inventory of inputs and constitutive principles}
\label{sec:inputs}
% ======================================================================

\subsection{PLEC components imported from Paper~1}

Paper~3 imports the four structurally necessary components of PLEC from
Paper~1 v4.0-PLEC~\cite{Paper1}.  They are used throughout this Supplement
in the following roles:

\begin{tcolorbox}[apfbox, title={Inputs, Constitutive Principles, and Honest Flags}]
\textbf{PLEC components (imported).}
\begin{itemize}
  \item \textbf{A1} (finite enforcement capacity, upper bound): $\sum_{d \in G}\varepsilon(d) \le C(\rho,\Gamma)$.  Used in every section.
  \item \textbf{MD} (positive cost floor, lower bound): $\inf_{d \in G}\varepsilon(d) \ge \mu^* > 0$.  Used in $T_\kappa$, $T_{\mathrm{entropy}}$.
  \item \textbf{A2} (argmin selection): $G_{\mathrm{realized}} = \argmin_{G \in \mathcal{A}} \sum \varepsilon$.  Used in entropy monotonicity (second law) and Fisher-metric uniqueness.
  \item \textbf{BW} (non-degeneracy): cost spectrum rich enough for $\argmin$ to be informative.  Used in $T_{\mathrm{seq}}$ temporal bound.
\end{itemize}

\textbf{Structural lemmas (imported).}
$T_{\mathrm{tensor}}$, $L_{\mathrm{loc}}$, $L_{\mathrm{nc}}$, $L_\Delta$,
$T_{\mathrm{adj}}$ (all from Paper~1); $L_{\mathrm{count}}$ ($C_{\mathrm{total}}=61$,
$d_{\mathrm{eff}}=102$), $T_{7B}$ (gauge-matter sector parity), $T_{\mathrm{holonomy}}$,
$B1'$ (real/pseudoreal composite structure) (all from Paper~2).

\textbf{Additional hypotheses introduced in this paper} (flagged at point of use):
\begin{itemize}
  \item \textbf{H1} --- Markovian-composition closure.  Composition of admissibility-preserving CPTP maps remains admissible.  Used in $T_{\mathrm{seq}}$ (\S\ref{sec:TCPTP}).
  \item \textbf{H2} --- Hamiltonian-constraint redundancy.  At a causal horizon, the Hamiltonian constraint identifies a gauge-redundant factor of $1/2$ in the ACC-to-entropy conversion.  Used in the $1/2 \times 1/2$ proof (\S\ref{sec:ACC}).
  \item \textbf{H3} --- Constraint-functional monotonicity.  Along admissible trajectories, $\Phi_c = \ln(d_{\mathrm{field}}/d_{\mathrm{admissible}})$ is monotone non-increasing.  Used in the staged-cosmogenesis arrow (\S\ref{sec:irr} closing).
\end{itemize}

None of H1--H3 is a PLEC component.  Each is a structural hypothesis local
to Paper~3; failures of any would not dissolve the framework but would
isolate the corresponding section.
\end{tcolorbox}

\subsection{What is assumed vs.\ what is proved}

\begin{center}
\small
\begin{tabular}{p{0.3\linewidth}p{0.3\linewidth}p{0.3\linewidth}}
\toprule
\textbf{Assumed (imported)} & \textbf{Proved in Paper~3} & \textbf{Deferred}\\
\midrule
A1, MD, A2, BW (PLEC) & $T_{\mathrm{CPTP}}$ & $\Lambda_{\mathrm{QCD}}$ (Paper~2)\\
$T_{\mathrm{tensor}}$, $L_{\mathrm{loc}}$ (Paper~1) & $T_\kappa$ ($\kappa=2$) & PMNS numerics (Paper~4)\\
$L_{\mathrm{nc}}$, $L_\Delta$, $T_{\mathrm{adj}}$ (Paper~1) & $T_{\mathrm{entropy}}$ & Quantitative cosmogenesis (Paper~6)\\
$L_{\mathrm{count}}$, $T_{7B}$ (Paper~2) & $T_{\mathrm{second\,law}}$ & Higher-order horizon corrections (Paper~6/7)\\
H1 (Markov closure) & $T_{\mathrm{seq}}$ (temporal Tsirelson) & Planck-scale closure (Paper~5)\\
H2 (Hamiltonian redundancy) & Horizon $S = A/4$ & Expansion history (Paper~6)\\
H3 ($\Phi_c$ monotonicity) & Fisher metric uniqueness & Matter-antimatter asymmetry (Paper~7)\\
\bottomrule
\end{tabular}
\end{center}

The ``proved'' column is what this Supplement establishes.  The ``assumed''
column is audited in Papers~1 and~2 (and in the Paper~1 essentiality
proofs~\cite{Paper1} for the four PLEC components).  The ``deferred''
column is explicitly out of scope for Paper~3.


% ======================================================================
\section{$T_{\mathrm{CPTP}}$: complete positivity as a theorem}
\label{sec:TCPTP}
% ======================================================================

The main paper (\S 4) presents the three-step proof; this
section fills in the proofs that were compressed and compares the result
explicitly with the Stinespring--Choi--GKSL infrastructure.

\subsection{Theorem statement}

\begin{tcolorbox}[apfbox, title={Theorem 2.1 ($T_{\mathrm{CPTP}}$)}]
Let $\Phi: \mathcal{B}(\mathcal{H}_S) \to \mathcal{B}(\mathcal{H}_S)$ preserve
admissibility of the enforcement regime at interface $\Gamma_S$.  Then $\Phi$
is completely positive and trace-preserving (CPTP), and admits a Kraus
decomposition $\Phi(\rho) = \sum_k K_k\rho K_k^\dagger$ with
$\sum_k K_k^\dagger K_k = \mathbb{1}$.
\end{tcolorbox}

\subsection{Full proof}

\textit{Step 1 --- Trace preservation.}  A1 fixes $C(\Gamma_S)$ as a finite
positive constant.  The normalization $\mathrm{tr}\,\rho = 1$ is
$C_{\mathrm{active}}/C(\Gamma_S) = 1$.  Suppose $\mathrm{tr}\,\Phi(\rho) > 1$.
Then committed capacity exceeds $C(\Gamma_S)$, violating A1's upper bound.
Suppose $\mathrm{tr}\,\Phi(\rho) < 1$.  Then committed capacity falls below
the fixed value; A1 requires both an upper bound and conservation of the
budget (because the budget is a constant of the interface, not merely a
ceiling).  Hence $\mathrm{tr}\,\Phi(\rho) = 1$.

\textit{Step 2 --- Complete positivity via ancilla-stability.}  By
$T_{\mathrm{tensor}}$ (Paper~1)~\cite{Paper1}, for any ancilla $R$ the
joint system $S \otimes R$ carries an independent budget $C(\Gamma_R)$.
The state $\rho_{SR} \geq 0$ is admissible iff all costs are non-negative.
Consider $\Phi \otimes \mathbb{1}_R$.  At $\Gamma_R$, the identity map
acts; costs are unchanged; admissibility at $\Gamma_R$ is preserved
trivially.  At $\Gamma_S$, $\Phi$ preserves admissibility by hypothesis.
By $L_{\mathrm{loc}}$ (Paper~1), the two interfaces are structurally
independent: evolution at one cannot create inadmissible states at the
other.  Hence the joint map is positive on $\rho_{SR} \geq 0$ --- i.e.,
$\Phi$ is completely positive.

\textit{Step 3 --- Kraus representation.}  Complete positivity + trace
preservation on a finite-dimensional Hilbert space (bounded by Paper~1's
$T_{\mathrm{form}}$/OR3) invokes Stinespring's dilation~\cite{Stinespring1955}:
there exists $\mathcal{H}_E$ and an isometry
$V: \mathcal{H}_S \to \mathcal{H}_S \otimes \mathcal{H}_E$ with
$\Phi(\rho) = \mathrm{tr}_E[V\rho V^\dagger]$.  Choi's
theorem~\cite{Choi1975} gives the equivalent Kraus decomposition
$\Phi(\rho) = \sum_k K_k \rho K_k^\dagger$ with $K_k = (\mathbb{1}_S \otimes \langle k|_E) V$.
The trace-preserving condition translates to $\sum_k K_k^\dagger K_k = \mathbb{1}$.
\hfill$\square$

\subsection{Comparison to Stinespring--Choi--GKSL}

A reviewer may reasonably ask: ``what does $T_{\mathrm{CPTP}}$ add to the
already-established Stinespring--Choi--GKSL theory?''  The answer is three
specific things.

\begin{enumerate}
  \item \textbf{Derivation direction.}  In conventional quantum information,
  the CPTP condition is \emph{imposed} as the definition of a valid
  evolution map, and Stinespring's theorem~\cite{Stinespring1955} then
  shows that every CPTP map admits a dilation.  In APF, the CPTP condition
  is \emph{derived} from admissibility preservation under A1 and
  $T_{\mathrm{tensor}}$, and Stinespring's theorem then supplies the Kraus
  form.  The direction of the derivation is reversed.
  \item \textbf{Physical interpretation of the dilation.}  Stinespring's
  environment $\mathcal{H}_E$ is traditionally treated as a mathematical
  device.  In APF, $\mathcal{H}_E$ is a second enforcement interface with
  its own capacity budget $C(\Gamma_E)$; the partial trace corresponds to
  abandoning the $\Gamma_E$ ledger.  The unitary evolution on
  $\mathcal{H}_S \otimes \mathcal{H}_E$ is the admissibility-preserving
  joint dynamics of two coupled budgets.
  \item \textbf{Restriction to finite dimension.}  Paper~1's $T_{\mathrm{form}}$
  bounds $\dim \mathcal{H}_S \le \lfloor C/\mu^* \rfloor$ at each interface.
  Stinespring's theorem applies in full generality (type $I_\infty$ and
  beyond); APF restricts to type $I_n$ with $n < \infty$.  The
  GKSL--Lindblad generators~\cite{Lindblad1976,GKS1976} then live in
  $\mathcal{B}(\mathcal{H}_S)$ with explicit finite trace.
\end{enumerate}

\noindent These three differences are what makes ``CPTP as a theorem'' a
nontrivial claim rather than a change of language.  The underlying
mathematics is unchanged; the provenance is.

\subsection{Sequential composition and the temporal Tsirelson bound}

\begin{tcolorbox}[apfbox, title={Theorem 2.2 ($T_{\mathrm{seq}}$)}]
Under H1 (Markovian-composition closure), any $T$-step sequence of admissible
CPTP maps $\Phi_T \circ \cdots \circ \Phi_1$ is itself admissible.
Time-evolved dichotomic observables $a(t) = U_t^\dagger a U_t$ satisfy the
Cirel'son operator identity; the temporal CHSH correlator is bounded by
$2\sqrt{2}$.
\end{tcolorbox}

\textit{Proof.}  CPTP maps form a semigroup under composition; the
composition of trace-preserving maps is trace-preserving; the tensor
product of two Kraus decompositions is a Kraus decomposition for the
composed map.  This is standard~\cite{Lindblad1976,GKS1976}; H1 is the
requirement that admissibility \emph{preservation} commutes with this
composition (admissibility-preserving at each step implies
admissibility-preserving for the composite).  For the temporal CHSH bound:
the Cirel'son identity $\|[a(t_1), a(t_2)]\|_{\mathrm{op}} \leq 2$ uses
only the algebraic structure of self-adjoint unitaries
(via $T_{\mathrm{adj}}$, Paper~1) and is state-independent, giving
$|\omega(\hat{S}_T)| \leq 2\sqrt{2}$ for the operator
$\hat{S}_T = a(t_1)\otimes(b(t_1)+b(t_2)) + a(t_2)\otimes(b(t_1)-b(t_2))$.
\hfill$\square$

\coderef{check\_T\_seq}{core.py}


% ======================================================================
\section{$T_\kappa$: the directed-enforcement multiplier $\kappa = 2$}
\label{sec:Tkappa}
% ======================================================================

\subsection{Theorem statement}

\begin{tcolorbox}[apfbox, title={Theorem 3.1 ($T_\kappa$)}]
In the directed-enforcement regime (one interface committed to enforcing a
distinction against a second interface), the per-distinction enforcement
multiplier satisfies $\kappa = 2$ uniquely.  Values $\kappa < 2$ correspond
to erasable single-interface commitment (not directed); values $\kappa > 2$
require an additional independent correlation that is not cost on the
original distinction.
\end{tcolorbox}

\subsection{Derivation of $\kappa = 2$}

The multiplier $\kappa$ answers: how much capacity is committed per
distinction when enforcing a directed (irreversible) commitment?  The
minimum at each interface is $\mu^*$ (MD); the directed enforcement commits
at each of two interfaces (the distinction-holder and the
distinction-witness).

\textit{Upper bound $\kappa \le 2$.}  A directed distinction requires
commitment at two interfaces only: the distinction-holding interface
(committing $\mu^*$ to maintain the distinction) and the witnessing
interface (committing $\mu^*$ to record the distinction).  A third
interface commitment would be an independent correlation --- a separate
distinction in its own right, not a higher $\kappa$ on the original.  By
Paper~1 $L_{\mathrm{nc}}$, multiple simultaneous commitments do not
compress into a single distinction's cost; they are separate entries in
the ledger.  Hence $\kappa \le 2$.

\textit{Lower bound $\kappa \ge 2$.}  If $\kappa < 2$, the directed
distinction is enforceable at less than $2\mu^*$.  Two sub-cases: $\kappa = 0$
is excluded by MD (no zero-cost distinctions); $0 < \kappa < 2$ would place
the commitment below the floor of two interfaces, which by
$L_{\mathrm{irr}}$ contradicts the irreversibility of directed commitment
(single-interface commitment is erasable; directed commitment is not).
Hence $\kappa \ge 2$.

\textit{Uniqueness.}  Both bounds saturate at $\kappa = 2$.  \hfill$\square$

\coderef{check\_T\_kappa}{core.py}

\subsection{Numerical verification}

The codebase evaluates $\kappa$ on the 61-mode ledger and returns
$\kappa = 2$ with no approximation: the computation is purely combinatorial
over the $C_{\mathrm{total}} = 61$ capacity types and yields
$\kappa = 2$ identically (see \coderef{check\_T\_kappa}{core.py}).  This is
\emph{not} a numerical coincidence: the derivation is structural, and the
code merely corroborates what the proof requires.


% ======================================================================
\section{$T_{\mathrm{entropy}}$: von Neumann entropy as committed capacity}
\label{sec:Tentropy}
% ======================================================================

\subsection{Theorem statement}

\begin{tcolorbox}[apfbox, title={Theorem 4.1 ($T_{\mathrm{entropy}}$)}]
The von Neumann entropy $S(\rho) = -\mathrm{tr}(\rho\log\rho)$ equals
$\kappa^{-1}$ times the enforcement demand of the active correlations at
the system interface:
\[
S(\rho) = \frac{1}{\kappa}\,\mathcal{E}_{\mathrm{cross}}(\rho) = \frac{1}{2}\,\mathcal{E}_{\mathrm{cross}}(\rho),
\]
where $\mathcal{E}_{\mathrm{cross}}$ counts the cross-interface enforcement
cost (in nats) committed by $\rho$.  The standard entropy properties
(positivity, concavity, subadditivity, strong subadditivity, data-processing
inequality) follow from A1, $L_\Delta$, $T_{\mathrm{CPTP}}$, and Lieb--Ruskai~\cite{LiebRuskai1973}.
\end{tcolorbox}

\subsection{Derivation}

\textit{Step 1 (Identification).}  By Paper~1 $T_{\mathrm{adj}}$, the
cost functional on the operator algebra $\mathcal{A}$ induces an inner
product; the cost of a projector $P_i$ is $\varepsilon(P_i) = \mathrm{tr}(\rho P_i)$
for a state $\rho$.  For a rank-one projector $P_i = |\psi_i\rangle\langle\psi_i|$
in the eigenbasis of $\rho$ with eigenvalue $\lambda_i$, $\varepsilon(P_i) = \lambda_i$.
The total cross-interface enforcement demand is
$\mathcal{E}_{\mathrm{cross}}(\rho) = -2 \sum_i \lambda_i \log \lambda_i$
(the factor of $\kappa = 2$ comes from directed commitment at both the
system and the environment interfaces; see $T_\kappa$).  Hence
$S(\rho) = \mathcal{E}_{\mathrm{cross}}/2 = -\sum_i \lambda_i \log \lambda_i$,
which is the von Neumann entropy.

\textit{Step 2 (Standard properties from imports).}
\begin{itemize}[nosep]
  \item \textbf{Positivity.}  $S(\rho) \geq 0$ from $\lambda_i \in [0,1]$.
  \item \textbf{Concavity.}  Follows from the concavity of $-x\log x$; preserved by cost-functional linearity on the algebra.
  \item \textbf{Subadditivity.}  $S(\rho_{SE}) \leq S(\rho_S) + S(\rho_E)$ follows from $L_\Delta$ applied to co-located distinctions at $\Gamma_S \cup \Gamma_E$.
  \item \textbf{Strong subadditivity.}  $S(\rho_{ABC}) + S(\rho_B) \leq S(\rho_{AB}) + S(\rho_{BC})$ is Lieb--Ruskai~\cite{LiebRuskai1973}; it follows in APF from the admissibility of the Naimark dilation of the tripartite structure and $T_{\mathrm{CPTP}}$.
  \item \textbf{Data processing.}  Under any CPTP map $\Phi$, the relative entropy is non-increasing~\cite{Wehrl1978}.  In APF, this is $T_{\mathrm{CPTP}}$ plus $L_{\mathrm{irr}}$: committed cross-interface correlations cannot decrease under local operations.
\end{itemize}
\hfill$\square$

\coderef{check\_T\_entropy}{core.py}

\subsection{Comparison to Lindblad and Spohn}

The mainstream open-systems route to the second law runs through Lindblad's
CP-map entropy inequalities~\cite{Lindblad1976} and Spohn's
positive-entropy-production framework for quantum dynamical
semigroups~\cite{Spohn1978}.  Spohn established that the entropy production
$\sigma(\rho) := -\frac{d}{dt} D(\rho_t \| \rho_*)$ is non-negative and
convex along the semigroup flow, where $\rho_*$ is the stationary state and
$D$ is the relative entropy; this gives approach to equilibrium for broad
classes of open quantum dynamics.

The framework here is in alignment, not in competition.  Three differences
of presentation, none of mathematics:
\begin{enumerate}
  \item \textbf{Source of irreversibility.}  Spohn locates irreversibility
  in the semigroup structure of the open-system dynamics.  The framework
  here locates it in $L_{\mathrm{irr}}$: local unrecoverability of
  cross-interface correlations.  These are two different routes to the
  same monotonicity statement.
  \item \textbf{Normalization.}  Spohn's entropy production is naturally
  formulated relative to the stationary state $\rho_*$.  The framework
  here works directly with the von Neumann entropy of $\rho_t$, with the
  factor of $\kappa = 2$ supplying the directed-enforcement multiplier
  that the open-systems literature usually absorbs into the choice of
  generator.
  \item \textbf{Scope.}  Spohn's framework is deeply developed for
  Markovian semigroups; the admissibility framework's $T_{\mathrm{seq}}$
  (\S\ref{sec:TCPTP}) extends to multi-step composition under hypothesis
  H1 (Markovian-composition closure), explicitly flagged as a separable
  hypothesis whose failure would isolate the multi-step second-law
  statement without affecting the single-step one.
\end{enumerate}
Paper~3's contribution to entropy theory is rhetorical and structural,
not a new monotonicity inequality.  The Lieb--Ruskai backbone is unchanged;
Spohn's entropy-production positivity is unchanged.  What changes is the
account of \emph{why} the monotonicity holds.

\subsection{The ledger-accounting principle (formal statement)}

\begin{tcolorbox}[apfbox, title={Proposition 4.2 (Ledger-Accounting Principle)}]
The entropy count is invariant under small changes in the enforcement-cost
landscape.  Formally: let $\varepsilon'(d) = \varepsilon(d) + \delta(d)$ be
a perturbation of the cost functional with $\|\delta\|_\infty < \mu^*$.
Then the count of admissible configurations at fixed capacity $C(\Gamma)$ is
unchanged: the admissible set under $\varepsilon'$ differs from that under
$\varepsilon$ only at the boundary, and the entropy functional $S(\rho) = \ln N_{\mathrm{admissible}}$
is invariant up to $O(\delta/\mu^*)$ corrections that vanish as the
perturbation vanishes.
\end{tcolorbox}

\textit{Sketch.}  The admissible set $\mathcal{A}$ is defined by two
inequalities: $\inf \varepsilon(d) \ge \mu^*$ and $\sum \varepsilon(d) \le C$.
Perturbations with $\|\delta\|_\infty < \mu^*$ preserve the floor
condition strictly; the budget condition shifts by $\sum \delta(d) = O(|G|\|\delta\|_\infty)$,
which for bounded cost spectra is bounded.  The count $|\mathcal{A}|$
is invariant up to boundary fluctuations of this order.  Entropy as
$\ln |\mathcal{A}|$ inherits the invariance.  \hfill$\square$

\noindent\textit{Remark.}  This is the formal statement of the claim
used in the main paper (\S 6.3, ledger-accounting principle) that entropy
is robust under cost-landscape corrections.  Paper~7's partition-function
treatment~\cite{Paper7} elevates this to a full thermodynamic argument;
here it functions as a direct count-based observation.


% ======================================================================
\section{$T_{\mathrm{second\,law}}$ and the arrow of time}
\label{sec:irr}
% ======================================================================

\subsection{$L_{\mathrm{irr}}$ (formal statement and proof)}

\begin{tcolorbox}[apfbox, title={Lemma 5.1 ($L_{\mathrm{irr}}$)}]
Let $S$ and $E$ be systems at independent interfaces $\Gamma_S$ and
$\Gamma_E$ (by $L_{\mathrm{loc}}$).  After an interaction producing
cross-interface correlation $\Delta(S,E) > 0$ (by $L_\Delta$), the
pre-interaction reduced state at $\Gamma_S$ is unrecoverable by any
CPTP map local to $\Gamma_S$.
\end{tcolorbox}

\textit{Proof.}  Suppose for contradiction that $\Phi_S$ is a
$\Gamma_S$-local CPTP map reversing the interaction as seen from
$\Gamma_S$: $\Phi_S(\rho_S^{\mathrm{after}}) = \rho_S^{\mathrm{before}}$.
The cross-interface capacity committed by the correlation is $\Delta > 0$;
this capacity is committed simultaneously at $\Gamma_S$ and $\Gamma_E$
(by $L_\Delta$).  For $\Phi_S$ to reclaim the full pre-interaction state,
it would need access to the $\Gamma_E$-side commitment.  By
$L_{\mathrm{loc}}$, $\Phi_S$ has no such access.  The recovered reduced
state at $\Gamma_S$ therefore differs from $\rho_S^{\mathrm{before}}$ by
the $\Gamma_E$-side capacity commitment, contradiction.  \hfill$\square$

\coderef{check\_L\_irr}{core.py}

\subsection{$T_{\mathrm{second\,law}}$ (formal statement)}

\begin{tcolorbox}[apfbox, title={Theorem 5.2 ($T_{\mathrm{second\,law}}$)}]
Along any admissible sequence of CPTP transitions
$\rho_0 \to \rho_1 \to \cdots \to \rho_T$, the von Neumann entropy is
non-decreasing:
\[
S(\rho_0) \le S(\rho_1) \le \cdots \le S(\rho_T).
\]
Equality $S(\rho_t) = S(\rho_{t+1})$ holds only when the $t$-th transition
commits no new cross-interface correlation.
\end{tcolorbox}

\textit{Proof.}  At each step, the transition $\rho_t \to \rho_{t+1}$ is
either (a) unitary (admits an inverse CPTP; equality case), or (b) involves
cross-interface coupling ($\Delta > 0$).  In case (b), $L_{\mathrm{irr}}$
guarantees the pre-transition state is unrecoverable; the data-processing
inequality~\cite{LiebRuskai1973,Wehrl1978} on the coupled system
$\rho_{SE,t}$ then forces $S(\rho_{S,t+1}) \geq S(\rho_{S,t})$.  Summing
over steps gives the chain.  \hfill$\square$

\coderef{check\_T\_second\_law}{core.py}

\subsection{The arrow of time and $T$-violation at $\pi/4$}

\begin{tcolorbox}[apfbox, title={Theorem 5.3 ($T_{\mathrm{CPT}}$)}]
The structural $T$-violation phase is $\theta_T = \pi/4$.  This value
follows from $L_{\mathrm{irr}}$, the holonomy-phase lemma
$L_{\mathrm{holonomy\,phase}}$ (Paper~2), and the algebraic structure of
Paper~1's T2 (finite-dimensional complex Hilbert-space representation).
\end{tcolorbox}

\textit{Sketch.}  The holonomy argument of Paper~2~\cite{Paper2} shows
that directed admissible cycles in the state space pick up phases
quantized in units related to the $\mathbb{Z}_2 \times \mathbb{Z}_2$
structure of the gauge sector.  The minimum directed-cycle $T$-violation
phase is half of $\pi/2$, i.e.\ $\pi/4$ --- the result of one $L_{\mathrm{irr}}$
step integrated around a smallest admissible directed loop.  \hfill$\square$

\coderef{check\_T\_CPT}{core.py}

\subsection{Macroscopic arrow from loss of global coordination}

\begin{tcolorbox}[apfbox, title={Proposition 5.4 (Macroscopic arrow)}]
At cosmological scale, a globally coordinated admissible state is one in
which admissibility holds jointly across all interfaces in a causal region.
Under H3 (monotonicity of the constraint functional $\Phi_c$), the set of
globally coordinated states is strictly decreasing along the admissible
trajectory: coordination can be lost but not regained.  The arrow of time
is the direction of coordination loss.
\end{tcolorbox}

\noindent This is $L_{\mathrm{irr}}$ applied to the coordinated-admissibility
observable over a macroscopic region; the proof is the same structure as
$L_{\mathrm{irr}}$, with ``local interface'' replaced by ``local causal
region.''


% ======================================================================
\section{Admissibility carrying capacity and $S = A/4$}
\label{sec:ACC}
% ======================================================================

\subsection{ACC: formal definition}

\begin{definition}[Admissibility Carrying Capacity]
Let $\Gamma$ be a relational interface (a causal boundary, a quantum
bipartition, or any co-dimension-1 surface across which admissibility is
enforced).  Its admissibility carrying capacity is
$C_{\mathrm{ACC}}(\Gamma) = $ the supremum number of distinct admissible
configurations that can coexist at the interface, weighted by their
cost in the enforcement currency.
\end{definition}

\subsection{The $1/2 \times 1/2$ bookkeeping}

\begin{tcolorbox}[apfbox, title={Theorem 6.1 (Horizon entropy prefactor)}]
At a causal horizon with area $A$, the entropy is
\[
S_H = \frac{1}{2} \cdot \frac{1}{2} \cdot \frac{A}{\ell_P^2} \cdot \frac{\ln d_{\mathrm{eff}}}{\ln 2}
= \frac{A}{4 \ell_P^2} \cdot \frac{\ln d_{\mathrm{eff}}}{\ln 2},
\]
which, after the conventional ACC-to-bit conversion, is $S_H = A/4$ in
natural units --- the Bekenstein--Hawking area law~\cite{Bekenstein1973,Hawking1975}.
\end{tcolorbox}

\textit{Proof.}  Two structurally forced factors of $1/2$ compose.

\textit{Factor 1 (Relational symmetrization).}  The ACC counts relational
distinctions across the interface.  Each distinction is a pair $(a, b)$
where $a$ is on one side and $b$ is on the other; under relational
equivalence $(a,b) \sim (b,a)$, each pair is counted twice in the naive
bookkeeping.  Symmetrization gives a factor of $1/2$.

\textit{Factor 2 (Hamiltonian-constraint redundancy; hypothesis H2).}
At a causal horizon, general covariance imposes the Hamiltonian
constraint: time reparameterizations at the horizon are gauge.  This
identifies two admissible configurations differing only in their gauge
slicing as physically identical, giving a second factor of $1/2$.  H2 is
the assumption that this identification is the \emph{only} relevant
gauge redundancy; other gauge symmetries (spatial diffeomorphisms,
internal gauge) are absorbed into the interface specification upstream.

Multiplying: $1/2 \times 1/2 = 1/4$, which is the Bekenstein--Hawking
prefactor.  \hfill$\square$

\coderef{check\_T\_deSitter\_entropy}{gravity.py}

\subsection{Unification of three entropies}

\begin{corollary}[Three-entropies unification]
Thermodynamic entropy, entanglement entropy, and horizon entropy are three
specializations of the same ACC-saturation count.  They differ only in
what the interface $\Gamma$ is:
\begin{itemize}
  \item Thermodynamic: $\Gamma$ is the system/environment boundary (implicit coarse-graining).
  \item Entanglement: $\Gamma$ is the quantum bipartition.
  \item Horizon: $\Gamma$ is the causal boundary.
\end{itemize}
\end{corollary}

\noindent\textit{Remark.}  This unification is a clear statement that the
framework subsumes rather than supplants standard black-hole and
cosmological-horizon thermodynamics.  The conventional $S = A/4$ law of
Bekenstein and Hawking~\cite{Bekenstein1973,Hawking1975} is recovered; the
de Sitter cosmological-horizon entropy of Gibbons and
Hawking~\cite{GibbonsHawking1977} is recovered as the saturated case
$S_{dS} = C_{\mathrm{total}} \ln d_{\mathrm{eff}} = 61 \ln 102$, with the
combinatorial values inherited from Paper~2's capacity-partitioning
argument.  The additional claim is that the same ledger underlies the
other two regimes.  The unification does not contradict any observation;
it reorganizes the existing explanations around a single counting
principle.

\subsection{Comparison to Bekenstein--Hawking and Gibbons--Hawking}

The standard derivations of horizon entropy operate at the semiclassical
level: Bekenstein's information-theoretic argument from the area
law~\cite{Bekenstein1973}, Hawking's particle-creation calculation
yielding the thermal spectrum at horizon temperature $T_H =
\hbar \kappa_{\mathrm{surf}} / 2\pi k_B$~\cite{Hawking1975}, and the
Gibbons--Hawking extension to cosmological event horizons in
positive-$\Lambda$ spacetimes with horizon temperature
$T_{dS} = H/2\pi$~\cite{GibbonsHawking1977}.  Each gives
$S = A/4$ in natural units.

The framework here is in the same lineage; the $1/2 \times 1/2$
decomposition is a specific combinatorial account of where the standard
$1/4$ prefactor comes from, valid under H2 (Hamiltonian-constraint
redundancy as the only relevant gauge structure in the horizon's
ACC-to-entropy conversion).  We do not provide a competing semiclassical
derivation, and the result is the same horizon entropy in either
language.  The framework's specific addition is the ACC vocabulary that
makes thermodynamic, entanglement, and horizon entropy three
specializations of one accounting (the corollary above); this is a
unification of three separate semiclassical results, not a re-derivation
of any of them.


% ======================================================================
\section{Fisher information and RG flow}
\label{sec:Fisher}
% ======================================================================

\subsection{Fisher metric uniqueness}

\begin{tcolorbox}[apfbox, title={Lemma 7.1 ($L_{\mathrm{Fisher}}$)}]
The unique Riemannian metric on the space of admissible density matrices
that (a) is monotone under CPTP maps, (b) reduces to the classical Fisher
metric on commuting states, and (c) is invariant under unitary conjugation,
is the quantum Fisher information metric (equivalently, the symmetric
logarithmic derivative / Bures metric) of Braunstein--Caves~\cite{BraunsteinCaves1994}.
Under the capacity-counting normalization of Paper~2 ($d_{\mathrm{eff}} = 102$),
\[
S = \frac{d_{\mathrm{eff}}}{2}\,\ln\det G,
\]
where $G$ is the Fisher metric on admissible states.
\end{tcolorbox}

\textit{Sketch.}  Monotonicity under CPTP pins the metric up to a one-parameter
family of admissible metrics in the Petz classification~\cite{Petz1996}.
Among these, the symmetric logarithmic derivative is distinguished by
(b)+(c), reducing to the classical Fisher metric on commuting states with
unitary invariance.  The entropy-determinant relation follows from the
capacity-counting structure of Paper~2~\cite{Paper2}: the entropy at
ACC saturation is $C_{\mathrm{total}} \ln d_{\mathrm{eff}}$, and the
volume form of $G$ scales as $\det G$ by standard arguments from
information geometry~\cite{AmariNagaoka2000}.  \hfill$\square$

\coderef{check\_L\_Fisher\_measure}{core.py}

\subsection{RG flow as Fisher gradient descent}

\begin{tcolorbox}[apfbox, title={Proposition 7.2 (RG flow as Fisher descent)}]
Under the Fisher metric $G$, the renormalization-group flow on the space
of admissible effective descriptions is the gradient descent of an entropy
functional.  Coarse-graining steps (integrating out high-momentum modes)
correspond to moves in the direction of maximum entropy gain per capacity
used --- the gradient of $S$ under $G$.
\end{tcolorbox}

\noindent\textit{Sketch.}  Coarse-graining is a CPTP map on the effective
theory.  By the data-processing inequality, entropy is non-decreasing under
coarse-graining.  Under the Fisher metric, the steepest-ascent direction for
entropy is the gradient; stationary points of the RG flow are fixed points
of this gradient descent.  The traditional RG beta-function interpretation
is recovered with the identification of the beta-function with the Fisher
gradient~\cite{AmariNagaoka2000}.  Full dynamical analysis is
Paper~4~\cite{Paper4} territory for specific flavor channels.

\subsection{Comparison to B\'eny--Osborne}

The closest external literature for the RG-as-Fisher-descent claim is the
program of B\'eny and Osborne~\cite{BenyOsborne2015}, who formulate RG
explicitly as a family of quantum channels with an induced information
geometry on the space of admissible effective descriptions.  In their
formulation, the RG step is a CPTP map, monotone RG-invariants are
information-theoretic quantities non-increasing along the flow, and the
Fisher metric provides the natural geometric object for characterizing
information loss.

The framework here is directionally aligned: RG as CPTP, Fisher metric as
the induced geometry, entropy gradient as the direction of RG flow.  The
additions specific to the admissibility framework are (i) the
capacity-count normalization ($d_{\mathrm{eff}} = 102$) from Paper~2's
structural argument and (ii) the connection to the Frobenius--Nagata
generation-channel structure of Paper~4.  Neither of these is provided
by the B\'eny--Osborne formulation; neither contradicts it.  Paper~3's
contribution is therefore a specific embedding of APF capacity counting
into the B\'eny--Osborne information-geometric RG picture, not a
re-derivation of that picture.


% ======================================================================
\section{Red-team section: additional hypotheses at point of use}
\label{sec:redteam}
% ======================================================================

Three hypotheses are used in this Supplement that are \emph{not} PLEC
components.  Each is flagged below at the point where it enters, with
the scope of what it supports and what would survive its failure.

\begin{tcolorbox}[redteambox, title={H1 --- Markovian-composition closure (\S\ref{sec:TCPTP})}]
\textbf{Statement.}  The composition of admissibility-preserving CPTP maps
is admissibility-preserving.

\textbf{Scope.}  Used in $T_{\mathrm{seq}}$ (temporal CHSH bound) and in the
second-law chain.  Without H1, the second law still holds step-by-step but
not as a global statement about multi-step evolutions.

\textbf{Status.}  [P] in Paper~7's partition-function framework; [C] in
Paper~3 alone.  The admissibility class is closed under composition if
and only if the composed map is itself admissibility-preserving.  Cases
where non-Markovian memory could violate this (long-range correlations at
a saturating interface) would falsify H1 for specific regimes.

\textbf{Red-team attack.}  Construct a pair of admissibility-preserving
maps whose composition is not admissibility-preserving.  None is known
within the framework, but the construction is open as a check on H1.
\end{tcolorbox}

\begin{tcolorbox}[redteambox, title={H2 --- Hamiltonian-constraint redundancy (\S\ref{sec:ACC})}]
\textbf{Statement.}  At a causal horizon, the only relevant gauge
redundancy in the ACC-to-entropy conversion is time reparameterization
(Hamiltonian constraint).  Other gauge symmetries are absorbed upstream.

\textbf{Scope.}  Provides the second factor of $1/2$ in the $1/2 \times 1/2$
bookkeeping for horizon entropy.  Without H2, the horizon-entropy
prefactor would require a different account; the ACC structure itself
would survive.

\textbf{Status.}  Consistent with general-relativistic canonical analysis
(ADM formulation).  [C] pending explicit derivation from Paper~6's
geometric admissibility framework.

\textbf{Red-team attack.}  Show a causal horizon where spatial
diffeomorphisms contribute an additional gauge redundancy beyond the
interface specification, changing the prefactor.  Plausible candidates:
rotating black holes (Kerr) with non-trivial angular gauge structure;
anti-de Sitter horizons with boundary conditions that modify the
counting.  Neither is worked out here.
\end{tcolorbox}

\begin{tcolorbox}[redteambox, title={H3 --- Constraint-functional monotonicity (\S\ref{sec:irr} closing)}]
\textbf{Statement.}  Along admissible trajectories, the constraint
functional $\Phi_c = \ln(d_{\mathrm{field}}/d_{\mathrm{admissible}})$ is
monotone non-increasing.

\textbf{Scope.}  Underwrites the macroscopic arrow of time and the staged
cosmogenesis picture of the main paper (\S 9 of the main paper).
Without H3, the quantum-interface arrow ($L_{\mathrm{irr}}$) still
survives; the cosmological-scale arrow would require an alternative
argument.

\textbf{Status.}  [P+H3] in this paper; [P] under the Paper~6 development
of the cost-landscape dynamics for gravitational degrees of freedom.

\textbf{Red-team attack.}  Exhibit an admissible cosmological trajectory
where $\Phi_c$ fails monotonicity --- e.g., a bouncing-cosmology model
that recovers earlier-epoch admissible configurations.  The framework
predicts no such trajectory exists.
\end{tcolorbox}


% ======================================================================
\section*{Appendix A: Theorem Index}
\addcontentsline{toc}{section}{Appendix A: Theorem Index}
\label{sec:supp_index}
% ======================================================================

\begin{center}
\small
\begin{tabular}{llll}
\toprule
\textbf{Result} & \textbf{Type} & \textbf{Depends on} & \textbf{Paper~3 section}\\
\midrule
$T_{\mathrm{CPTP}}$ & Theorem & A1, $T_{\mathrm{tensor}}$, $L_{\mathrm{loc}}$ & \S\ref{sec:TCPTP}\\
$T_{\mathrm{seq}}$ & Theorem & $T_{\mathrm{CPTP}}$, $T_{\mathrm{adj}}$, H1 & \S\ref{sec:TCPTP}\\
$T_\kappa$ & Theorem & MD, $L_{\mathrm{irr}}$, $L_{\mathrm{nc}}$ & \S\ref{sec:Tkappa}\\
$T_{\mathrm{entropy}}$ & Theorem & $T_\kappa$, $T_{\mathrm{adj}}$, $L_\Delta$, Lieb--Ruskai~\cite{LiebRuskai1973} & \S\ref{sec:Tentropy}\\
Ledger-accounting Prop. & Proposition & A1, MD & \S\ref{sec:Tentropy}\\
$L_{\mathrm{irr}}$ & Lemma & A1, $L_{\mathrm{nc}}$, $L_{\mathrm{loc}}$, $L_\Delta$ & \S\ref{sec:irr}\\
$T_{\mathrm{second\,law}}$ & Theorem & $L_{\mathrm{irr}}$, $T_{\mathrm{entropy}}$ & \S\ref{sec:irr}\\
$T_{\mathrm{CPT}}$ ($\theta_T = \pi/4$) & Theorem & $L_{\mathrm{irr}}$, $T_{\mathrm{holonomy}}$, T2 & \S\ref{sec:irr}\\
Macroscopic arrow & Proposition & $L_{\mathrm{irr}}$, H3 & \S\ref{sec:irr}\\
Horizon $S = A/4$ & Theorem & ACC, H2, Bekenstein~\cite{Bekenstein1973}, Hawking~\cite{Hawking1975} & \S\ref{sec:ACC}\\
Three-entropies unification & Corollary & $T_{\mathrm{entropy}}$, ACC & \S\ref{sec:ACC}\\
$L_{\mathrm{Fisher}}$ & Lemma & Braunstein--Caves~\cite{BraunsteinCaves1994}, Petz~\cite{Petz1996}, $T_{\mathrm{CPTP}}$ & \S\ref{sec:Fisher}\\
RG flow as Fisher descent & Proposition & $L_{\mathrm{Fisher}}$, data processing & \S\ref{sec:Fisher}\\
\bottomrule
\end{tabular}
\end{center}


% ======================================================================
\section*{Appendix B: Dependency diagram}
\addcontentsline{toc}{section}{Appendix B: Dependency diagram}
\label{sec:supp_deps}
% ======================================================================

\begin{center}
\small
\begin{verbatim}
                         A1 + MD + A2 + BW  (PLEC, from Paper 1)
                                 |
                    +------------+------------+
                    |                         |
            T_tensor (Paper 1)         L_nc, L_Delta, L_loc (Paper 1)
                    |                         |
                    v                         v
               T_CPTP            L_irr  +---> T_kappa (kappa = 2)
                    |              |                |
                    | H1           |                v
                    v              |        T_entropy <- Lieb-Ruskai [LR73]
                  T_seq            |                |
                                   v                v
                            T_second_law    Ledger-accounting
                                   |
                           +-------+-------+
                           |               |
                      T_CPT (pi/4)    Macroscopic arrow (H3)
                                        |
                                        v
                               Cosmogenesis (Paper 3 main)

           ACC + H2 -> Horizon S = A/4 -> Three-entropies unification
                  |
                  |
                  +-- L_Fisher -> RG flow as Fisher descent
                        |
                        v
                   (Paper 4 flavor channels)
\end{verbatim}
\end{center}


% ======================================================================
\begin{thebibliography}{9}
% ======================================================================

\bibitem{Stinespring1955}
W.\ F.\ Stinespring, ``Positive functions on $C^*$-algebras,''
\emph{Proceedings of the American Mathematical Society} \textbf{6}, 211--216 (1955).

\bibitem{Choi1975}
M.-D.\ Choi, ``Completely positive linear maps on complex matrices,''
\emph{Linear Algebra and its Applications} \textbf{10}, 285--290 (1975).

\bibitem{Lindblad1976}
G.\ Lindblad, ``On the generators of quantum dynamical semigroups,''
\emph{Communications in Mathematical Physics} \textbf{48}, 119--130 (1976).

\bibitem{GKS1976}
V.\ Gorini, A.\ Kossakowski, and E.\ C.\ G.\ Sudarshan, ``Completely positive dynamical semigroups of $N$-level systems,''
\emph{Journal of Mathematical Physics} \textbf{17}, 821--825 (1976).

\bibitem{Kraus1983}
K.\ Kraus, \emph{States, Effects, and Operations}, Springer, 1983.

\bibitem{LiebRuskai1973}
E.\ H.\ Lieb and M.\ B.\ Ruskai, ``Proof of the strong subadditivity of quantum-mechanical entropy,''
\emph{Journal of Mathematical Physics} \textbf{14}, 1938--1941 (1973).

\bibitem{Spohn1978}
H.\ Spohn, ``Entropy production for quantum dynamical semigroups,''
\emph{Journal of Mathematical Physics} \textbf{19}, 1227--1230 (1978).

\bibitem{Wehrl1978}
A.\ Wehrl, ``General properties of entropy,''
\emph{Reviews of Modern Physics} \textbf{50}, 221--260 (1978).

\bibitem{Bekenstein1973}
J.\ D.\ Bekenstein, ``Black holes and entropy,''
\emph{Physical Review D} \textbf{7}, 2333--2346 (1973).

\bibitem{Hawking1975}
S.\ W.\ Hawking, ``Particle creation by black holes,''
\emph{Communications in Mathematical Physics} \textbf{43}, 199--220 (1975).

\bibitem{GibbonsHawking1977}
G.\ W.\ Gibbons and S.\ W.\ Hawking, ``Cosmological event horizons, thermodynamics, and particle creation,''
\emph{Physical Review D} \textbf{15}, 2738--2751 (1977).

\bibitem{BraunsteinCaves1994}
S.\ L.\ Braunstein and C.\ M.\ Caves, ``Statistical distance and the geometry of quantum states,''
\emph{Physical Review Letters} \textbf{72}, 3439--3443 (1994).

\bibitem{Petz1996}
D.\ Petz, ``Monotone metrics on matrix spaces,''
\emph{Linear Algebra and its Applications} \textbf{244}, 81--96 (1996).

\bibitem{AmariNagaoka2000}
S.-I.\ Amari and H.\ Nagaoka, \emph{Methods of Information Geometry}, AMS/Oxford, 2000.

\bibitem{BenyOsborne2015}
C.\ B\'eny and T.\ J.\ Osborne, ``The renormalization group via statistical inference,''
\emph{New Journal of Physics} \textbf{17}, 083005 (2015).

\bibitem{Paper1}
E.\ S.\ Brooke, ``The Enforceability of Distinction,''
v4.0-PLEC (2026). \url{https://doi.org/10.5281/zenodo.18604678}

\bibitem{Paper2}
E.\ S.\ Brooke, ``The Structure of Admissible Physics,''
v5.3-PLEC (2026). \url{https://doi.org/10.5281/zenodo.18604839}

\bibitem{Paper4}
E.\ S.\ Brooke, ``Admissibility Constraints and Structural Saturation,''
Zenodo deposit (2026). \url{https://doi.org/10.5281/zenodo.18604845}

\bibitem{Paper7}
E.\ S.\ Brooke, ``Action, Internalization, and the Lagrangian,''
v2.0 (2026). \url{https://doi.org/10.5281/zenodo.18604875}

\end{thebibliography}

\end{document}
